\subsection{The checked expansion invariant and a sharper obstruction}

\begin{proposition}[Known, with source corrections: the repaired AKR invariant]
\label{C:akr-invariant}
Let $K^*=\bigcup_mK_m$ be the fresh-target construction of
\citet[Construction C.1]{AdamekKoubekReiterman1981}, with the
actual-path repair in TOTEX Proposition C.1, and any admissible
sequence $a_0=2>a_1>\cdots>3/2$ and quasi-filter parameters.
For an object $X\in K_m$, an arrow $g:Y\to X$, and a finite
family of pairwise distinct arrows $h_i:W\to X$, we have:
\begin{enumerate}[label=(\alph*)]
\item if $|\im h\cap\im g|\geq2$ for an arrow $h:W\to X$,
then $h=gh'$ for an arrow $h':W\to Y$;
\item if $|\im h_i\cap\im g|\leq1$ for every $i$, then
$|\bigcup_i\im h_i\setminus\im g|\geq a_m r$;
\item for every nonempty such family,
$|\bigcup_i\im h_i|>3r/2$.
\end{enumerate}
The maps in (a) and the equality $h=gh'$ are maps of the
concrete category of finite sets and injections, not merely
factorizations of sets after forgetting the category.
\end{proposition}
\begin{proof}[Verification of the cited induction]
We opened the full proof of Theorem C.2 on pp.~385--389,
including the induction on $(c_m),(d_m)$ that was not completed
in the previous review. These are precisely (a),(b). We record
the operative steps so that the use of the theorem is explicit.

At the initial stage the only nonidentity arrows are the three
points of $W$, so (a),(b) hold with $a_0=2$. For a new
quasi-filter target $Z$, every old arrow into $Z$ factors through
one of its copies $\alpha_i:T\to Z$. If a $W$-map meets a
specified copy in at least two points, overlap condition (10)
allows us to move it one neighboring copy toward that copy.
Condition (11) identifies the neighboring overlap with the
image of the relevant map $S\to T$. Applying the old (a) to
that overlap gives a morphism $W\to S$, and hence the required
transported $W$-map in the neighboring copy. Repetition moves
the whole arrow into the specified copy. The old (a) there
proves the new (a), including the categorical factorization.

For (b), assign each selected $W$-map to its first possible
copy and partition the cyclic list of $s$ copies into blocks
of $3n$. Remove a block carrying the fewest assigned arrows;
at most $r\,3n/s$ arrows are lost. Traverse the remaining copies
from the two ends toward the removed block. At each step,
(10) puts all intersection with the previously visited copies
inside the neighboring overlap. A newly assigned $W$-map
cannot meet that overlap in two points: (a), (11), and
monicity would make it factor through the preceding copy,
contrary to its assignment. The old (b) therefore contributes
at least $a_m$ new points for every retained arrow. Hence
\[
 \left|\bigcup_i\im h_i\setminus\im g\right|
 \geq a_mr\left(1-\frac{3n}{s}\right)
 \geq a_{m+1}r,
\]
where the last inequality is exactly the parameter restriction
in Construction C.1. The two traversal ranges and their
indices were checked directly on the page images on
pp.~387--389. For a new disjoint sum, arrows factor through
one summand, and the two summands contribute disjoint sets;
the old bounds give the claimed bound immediately.

For (c), distinct $W$-maps meet in at most one point: otherwise
(a) factors one through the other, and the only endomorphism
of $W$ is its identity. Applying (b) with $g=h_1$ to the
remaining $r-1$ arrows yields
\[
 \left|\bigcup_{i=1}^{r}\im h_i\right|
 \geq 3+a_m(r-1)>\frac32r.
\]
This also justifies the union bound used in the disjoint-sum
case of the printed induction.

The actual-path repair supplies (9)--(11); we checked the
formula and the interval overlap verification in TOTEX
Proposition C.1. Two further index misprints in the source do
not affect the induction: on p.~387, the definition of $L$
must use $\{1,\ldots,r\}$, not $\{1,\ldots,n\}$; $r$
counts selected $W$-maps and $n$ is the unrelated path bound.
On p.~388, ``Then $(g_{m+1})$ follows'' must refer to
$(g_t)$. These corrections were confirmed against the page
images, not inferred only from OCR.
\end{proof}

\begin{proposition}[Proved: pointwise splitting of the AKR column]
\label{C:akr-split-column}
Let $R$ be a nonzero commutative unital ring, let
$P_x^R=R[K^*(x,-)]$, and define a morphism of covariant
$R$-module diagrams
\[
 d^T:P_W^R\longrightarrow(P_V^R)^2,\qquad
 t\longmapsto\bigl(t(v_2-v_1),t(v_3-v_1)\bigr).
\]
For every object $X$, $d^T_X$ is a split injection of finite
free $R$-modules, and its cokernel $Q(X)$ is finite free.
No natural splitting is asserted. For the diagram
$Q=\coker d^T$, derived colimits satisfy
\[
 \colim_{K^*}Q\cong R^2,\qquad
 L_1\colim_{K^*}Q\cong R,\qquad
 L_i\colim_{K^*}Q=0\quad(i\geq2).
\]
Thus the failure of exactness is witnessed by an objectwise
free diagram with a projective resolution of length one.
\end{proposition}
\begin{proof}
Every object $X$ is introduced at a finite stage. No later
stage adds arrows between old objects, so $K^*(W,X)$ and
$K^*(V,X)$ are finite. For a nonempty set $S$ of $r$ columns,
Proposition~\ref{C:akr-invariant}(c) says the corresponding
$3r$ incidences $gv_i$ use more than $3r/2$ points. Hence
some point belongs to exactly one of those incidences. If its
color is $v_2$ or $v_3$, the corresponding row in the first
or second coordinate has a sole entry $+1$ among these
columns. If its color is $v_1$, either corresponding row
has a sole entry $-1$.

Choose such a row, remove its unique column, and repeat.
The selected rows are distinct, because each previously
selected row is zero on all remaining columns. Ordering the
rows and columns in this elimination order gives a triangular
square minor with diagonal entries $\pm1$. It is invertible
over $\Z$, hence over every $R$. Projection onto the selected
rows followed by its inverse is a left inverse of $d_X^T$.
Eliminating these selected coordinates identifies the quotient
with a free module on the unselected rows. This proves
pointwise splitting and freeness without selecting splittings
natural in $X$.

Thus
\[
 0\longrightarrow P_W^R\xrightarrow{d^T}(P_V^R)^2
 \longrightarrow Q\longrightarrow0
\]
is an exact sequence of diagram modules, and its two left
terms are projective. For every object $x$, the augmentation
identifies $\colim P_x^R$ with $R$: the category $x/K^*$ has
initial object $\id_x$. Under these identifications the
induced map $R\to R^2$ is zero, since both coordinate
differences have augmentation zero. Applying colimit to
this length-one projective resolution gives exactly the
three displayed derived-colimit groups. This is compatible
with non-totality, but it does not prove non-totality: the
coefficient diagram $Q$ has not been shown to lie in the
acyclic class generated by permutation diagrams.
\end{proof}
